\documentclass[12pt,a4paper]{article}
\usepackage{fontspec}
\usepackage{polyglossia}
\setmainlanguage{russian}
\setmainfont{Liberation Serif}
\usepackage[left=0.9cm,right=0.9cm,top=0.8cm,bottom=0.8cm]{geometry}
\usepackage{anyfontsize}
\usepackage{amsmath}
\usepackage{array}
\setlength{\parindent}{0pt}
\pagestyle{empty}
\renewcommand{\arraystretch}{1.65}

% ТОЛЬКО НОМЕР И ОТВЕТ. Ничего больше сюда не добавлять:
% разборы, прогноз, уровни задач и пояснения живут в SESSIYA.md и prognoz.py.
% Лист читается с телефона одним взглядом — лишняя строка ломает эту цель.
% Внимание: в дополнительных (10–19) под одним номером в A и Б стоят РАЗНЫЕ задачи.

\begin{document}

\noindent{\fontsize{26}{30}\selectfont\bfseries Комбинаторика~1\ \ ответы}\hfill{\large 22 сентября, 7\,И}
\vspace{5pt}\hrule\vspace{9pt}
\fontsize{19}{19}\selectfont

\noindent
\begin{tabular}{@{}r@{\ \ }p{9.2cm}p{9.2cm}@{}}
\textbf{№} & \textbf{вариант A} & \textbf{вариант Б} \\[2pt]
\hline
 1 & 107 \quad 141 \quad \textbf{нельзя}      & 109 \quad 155 \quad \textbf{нельзя} \\
 2 & 33 \quad 67 \quad 53 \quad 53\,\%        & 37 \quad 113 \quad 75 \quad 50\,\% \\
 3 & 18                                        & 22 \\
 4 & 36 \quad 45                               & 45 \quad 15 \\
 5 & 127 \ 136 \ 145 \ 235                        & 129 \ 138 \ 147 \ 156 \ 237 \ 246 \ 345 \\
 6 & $n-m+1$                                   & $b-a+1$ \\
 7 & 1991                                      & 2971 \\
 8 & 118                                       & 104 \\
 9 & 6                                         & 9 \\[3pt]
\hline
10 & 11 \quad 40                               & 36 \quad 110 \\
11 & 4, 7, 9 \quad 10                          & 420 \quad 400 \\
12 & 40 \quad 600 \quad 542                    & 4 \quad 12 \quad 8 \\
13 & 3 \quad 10 \quad 7                        & 15 \quad 35 \\
14 & 22 \quad 84                               & 22 \quad 64 \\
15 & \textbf{нет} \quad не меньше 5            & 27 \quad 93 \\
16 & 0, 2 или 4                                & 60 \quad 900 \quad 813 \\
17 & 240 \quad 225                             & 4, 7, 10 \quad 7 \\
18 & 26 \quad 74                               & \textbf{нет} \quad не меньше 5 \\
19 & 10 \quad 15                               & 1, 3, 5 или 7 \\[3pt]
\hline
20 & \multicolumn{2}{@{}l@{}}{5 \ 6 \ 7 \ 8 \ 9 \ 10 \ 12 \ 13 \ 14 \ 15 \quad--- единственное} \\
21 & \multicolumn{2}{@{}l@{}}{6210001000 \quad--- единственное} \\
22 & \multicolumn{2}{@{}l@{}}{доказательство, принимаем устно} \\
\hline
\end{tabular}

\vspace{7pt}
\noindent{\fontsize{13}{15}\selectfont В задачах 10--19 под одним номером в вариантах A и Б стоят \emph{разные} задачи.}

\end{document}
